\documentclass[a0,portrait]{a0poster}

\usepackage[T1]{fontenc}
\usepackage{lmodern}

\usepackage[margin=4cm]{geometry}

\usepackage{mathtools,amssymb}

\usepackage{graphicx}
\usepackage{xcolor}
\usepackage{tikz}
\usetikzlibrary{calc, positioning}

\usepackage{titlesec}
\usepackage{booktabs}
\usepackage{caption}

\usepackage[hidelinks,pdfusetitle]{hyperref}
\usepackage[all]{hypcap}
\usepackage[capitalize]{cleveref}

\usepackage{multicol}
\usepackage[export]{adjustbox}
\usepackage[framemethod=TikZ]{mdframed}

\columnsep=32pt

\definecolor{custom_red}{HTML}{d81e2c}

\titleformat{\section}{\color{custom_red}\Large\bfseries}{}{24pt}{}

\mdfdefinestyle{MyFrame}{
    linecolor=custom_red,
    outerlinewidth=2pt,
    roundcorner=48pt,
    innertopmargin=\baselineskip,
    innerbottommargin=\baselineskip,
    innerrightmargin=24pt,
    innerleftmargin=24pt,
    backgroundcolor=white!50!white,
}

\begin{document}

\begin{mdframed}[style=MyFrame]

    \vspace{0.5cm}

    \hspace{0.5cm}
    \begin{minipage}[b]{0.4\linewidth}
        \includegraphics[height=4cm,valign=t]{images/logo-EPS.png}
    \end{minipage}
    \hfill
    \begin{minipage}[b]{0.4\linewidth}
        \raggedleft
        \includegraphics[height=4cm,valign=t]{images/logo-UCM.png}
    \end{minipage}
    \hspace{1cm}

    \vspace{2cm}

    \begin{minipage}[h]{0.98\linewidth}
        \centering \huge \textcolor{custom_red}{
            \textbf{Evolutionary adaptation through bet-hedging in finite populations under fluctuating environments}
        }\\
        \smallskip
        \Large \underline{Marco Mendívil-Carboni}$^1$, Juan J. Mazo$^1$ and Luis Dinis$^1$\\
        \smallskip
        \normalsize $^1$Grupo Interdisciplinar de Sistemas Complejos (GISC) and Departamento de Estructura de la Materia, Física Térmica y Electrónica, Universidad Complutense de Madrid, 28040 Madrid, Spain\\
    \end{minipage}

    \vspace{1.5cm}

    \large % Elevates overall font scaling within the body
    \begin{multicols}{3}

        \section{Introduction \& Objectives}
        \begin{itemize}
            \item Natural populations face unpredictable, fluctuating environments under finite resource constraints.
            \item While classic evolutionary theory suggests selection favors strategies maximizing mean growth rate $\langle \mu \rangle$, finite populations subject to demographic constraints face non-negligible extinction risks.
            \item \textbf{Key Question:} Does natural selection favor growth-maximizing strategies or fluctuation-reducing resilience strategies when population capacity is limited?
        \end{itemize}

        \section{Stochastic Model Framework}
        \begin{itemize}
            \item \textbf{Environment $e \in \{1,2\}$:} Continuous-time Markov process with transition rates $\omega_t(e,e')$.
            \item \textbf{Population ($N \le N_\text{ini}$):} Birth-death process where individuals carry strategy $s(A) = 1 - s(B)$ regulating phenotypic diversity.
            \item \textbf{Mutations:} Offspring inherit parent strategy or mutate with probability $p_\text{mut}$.
        \end{itemize}

        % \vspace{0.5em}
        % \begin{table}[H]
        %     \centering
        %     \caption*{\textbf{Persister--Resistant (PR) Parameters}}
        %     \small
        %     \begin{tabular}{c c c}
        %         \toprule
        %         Rates                    & Env 1 (Favorable)  & Env 2 (Harsh)      \\
        %         \midrule
        %         Birth $\omega_b(e,\phi)$ & $A: 1.0,\; B: 0.2$ & $A: 0.0,\; B: 0.0$ \\
        %         Death $\omega_d(e,\phi)$ & $A: 0.0,\; B: 0.0$ & $A: 1.0,\; B: 0.1$ \\
        %         \bottomrule
        %     \end{tabular}
        % \end{table}

        \section{Growth vs. Survival Trade-off}
        \begin{center}
            \includegraphics[scale=2.986]{figures/asymmetric/strat_phe_0_i/rates.pdf}
            \captionof{figure}{\textbf{Pareto Front:} Extinction rate $r_\text{ext}$ vs. growth rate $\langle\mu\rangle$. Evolved strategies select a middle-ground compromise between growth and survival.}
        \end{center}

        \columnbreak

        % ==================== COLUMN 2 ====================
        \section{Strategy Trajectories \& Fixation}
        \begin{itemize}
            \item \textbf{Rapid Fixation:} Advancing mutant lineages quickly dominate the population structure.
            \item \textbf{Strategy Drift:} Narrow strategy variance $\sigma_s$ leads to coherent exploration of strategy space over long time scales.
        \end{itemize}

        \begin{center}
            \includegraphics[scale=2.986]{figures/asymmetric/time_series/avg_strat_phe_0.pdf}
            \captionof{figure}{\textbf{Time Series Dynamics:} Mean strategy trajectory over time. Dotted vertical lines mark stochastic extinction events.}
        \end{center}

        \section{Population Constraint Modulation}
        \begin{itemize}
            \item \textbf{Compromise Strategy:} Evolution selects strategies strictly between growth maximization ($s(A) \sim 0.8$) and extinction minimization ($s(A) \sim 0.0$).
            \item \textbf{Scaling Transition:} Large $N_\text{ini}$ converges to growth maximization, whereas small $N_\text{ini}$ favors resilient strategies to avoid extinction.
        \end{itemize}

        \begin{center}
            \includegraphics[scale=2.986]{figures/asymmetric/n_agents_i/avg_avg_strat_phe_0.pdf}
            \captionof{figure}{\textbf{Transition:} Evolved strategy $\langle\bar{s}\rangle(A)$ vs. initial carrying capacity $N_\text{ini}$.}
        \end{center}

        \section{Heuristic Interpretation}
        \begin{itemize}
            \item Long-term strategy probability distribution $p(\bar{s})$ balance:
                  \begin{equation*}
                      p(\bar{s}) \propto \frac{\exp(N_\text{ini} \cdot \langle\mu\rangle(\bar{s}))}{r_\text{ext}(N_\text{ini},\bar{s}) + p_\text{mut} \cdot \langle\mu_b\rangle(\bar{s})}
                  \end{equation*}
            \item \textbf{Scaling Mismatch:} Growth selection enters \textbf{exponentially} in the numerator, while extinction rate scales as a \textbf{power-law} $r_\text{ext} \sim N_\text{ini}^{-\alpha}$ in the denominator.
            \item Extinction risk dominates at low $N_\text{ini}$; growth selection inevitably prevails as $N_\text{ini} \to \infty$.
        \end{itemize}

        \section{7. Robustness Analysis}
        \begin{itemize}
            \item \textbf{Incremental Mutations:} Continuous Gaussian mutations ($\sigma_\text{mut}=0.1$) yield qualitatively identical strategy transitions.
            \item \textbf{Symmetric Adaptation Set:} Re-evaluated under symmetric environments, confirming population constraints generalizable across different fitness landscapes.
        \end{itemize}

        \begin{center}
            \includegraphics[scale=2.986]{figures/symmetric/n_agents_i/avg_avg_strat_phe_0.pdf}
            \captionof{figure}{\textbf{Symmetric Adaptation:} Evolved strategy under the Symmetric-Adaptation (SA) parameter set.}
        \end{center}

        \section{8. Conclusions}
        \begin{itemize}
            \item Mean growth rate alone is insufficient to predict evolutionary outcomes in finite populations.
            \item Demographic constraints shift selection from growth maximization to fluctuation avoidance/resilience.
            \item The $N_\text{ini}$ transition is dictated by the scaling mismatch between exponential growth selection and power-law extinction risks.
        \end{itemize}

    \end{multicols}

    \vspace{1cm}

    \begin{center}
        \color{custom_red}{\large 3rd Biology for Physics \ -- \ Barcelona \ -- \ September 6 - 10, 2026}
    \end{center}

\end{mdframed}

\end{document}
